\documentclass[12pt]{article}
\usepackage{amsmath,amssymb,amsfonts}
\usepackage[margin=1in]{geometry}

\begin{document}

\section*{Chapter 3, Problem 1}

\textbf{Problem.} Prove that the complex $n$th roots of unity together with multiplication of complex numbers form an abelian group.

\textit{[Hint: From De Moivre's formula, the complex $n$th roots of unity are given by}
\[
\varepsilon_k = \cos\frac{2\pi k}{n} + i\sin\frac{2\pi k}{n} = e^{2\pi i k/n}, \quad k = 0, 1, \ldots, n-1,
\]
\textit{where $i^2 = -1$.]}

\bigskip
\textbf{Solution.}

Let $G = \{\varepsilon_0, \varepsilon_1, \ldots, \varepsilon_{n-1}\}$ where $\varepsilon_k = e^{2\pi i k/n}$. We verify the four group axioms under multiplication:

\medskip
\textbf{1. Closure:} For any $\varepsilon_j, \varepsilon_k \in G$,
\[
\varepsilon_j \cdot \varepsilon_k = e^{2\pi i j/n} \cdot e^{2\pi i k/n} = e^{2\pi i (j+k)/n} = e^{2\pi i [(j+k) \bmod n]/n} = \varepsilon_{(j+k) \bmod n}.
\]
Since $0 \le (j+k) \bmod n \le n-1$, we have $\varepsilon_j \cdot \varepsilon_k \in G$.

\medskip
\textbf{2. Associativity:} Multiplication of complex numbers is associative, so for any $\varepsilon_j, \varepsilon_k, \varepsilon_l \in G$:
\[
(\varepsilon_j \cdot \varepsilon_k) \cdot \varepsilon_l = \varepsilon_j \cdot (\varepsilon_k \cdot \varepsilon_l).
\]

\medskip
\textbf{3. Identity:} The element $\varepsilon_0 = e^{0} = 1$ satisfies
\[
\varepsilon_0 \cdot \varepsilon_k = \varepsilon_k \cdot \varepsilon_0 = \varepsilon_k
\]
for all $\varepsilon_k \in G$.

\medskip
\textbf{4. Inverse:} For each $\varepsilon_k \in G$, consider $\varepsilon_{n-k} \in G$:
\[
\varepsilon_k \cdot \varepsilon_{n-k} = e^{2\pi i k/n} \cdot e^{2\pi i (n-k)/n} = e^{2\pi i n/n} = e^{2\pi i} = 1 = \varepsilon_0.
\]
Thus $\varepsilon_k^{-1} = \varepsilon_{n-k} \in G$.

\medskip
\textbf{5. Commutativity (Abelian):} Since multiplication of complex numbers is commutative,
\[
\varepsilon_j \cdot \varepsilon_k = \varepsilon_k \cdot \varepsilon_j
\]
for all $\varepsilon_j, \varepsilon_k \in G$.

\medskip
Therefore, $(G, \cdot)$ is an abelian group. \hfill $\blacksquare$

\end{document}
